\scanpage{317}
\mspage{267}

on a

\[
\deleted{\unsure{\alpha}}g\in W\Longleftrightarrow\overline{\deleted{\unsure{Z}}g}\in W\quad\text{i.e.}\quad Z(g)\ W\text{-asph.},
\]

en particulier

\[
\deleted{\unsure{\alpha}}f_i\in W\Longleftrightarrow Z_i\ \deleted{\unsure{d}}W\text{-asph.}
\]

Ainsi, l'énoncé à prouver sur le cube initial se ramène à \deleted{\unsure{celle}} un énoncé \deleted{\uline{Corollaire} 1. Moyennant $W(7\ \mathrm{bis})$, si on a un diagramme} sur le carré

\diagram{\(\begin{tikzcd}[ampersand replacement=\&,row sep=small,column sep=small]
{}\&Z_1\arrow[dr]\&{}\\
Z_0\arrow[ur]\arrow[dr,dash]\&{}\&Z_3\\
{}\&Z_2\arrow[ur]\&{}
\end{tikzcd},\)}

savoir que les trois conditions

\[
\begin{aligned}
&\text{a}')\quad Z_0\ \deleted{\unsure{d}}W\text{-asph.}\\
&\text{b}')\quad Z_1\text{ et }Z_2\quad W\text{-asph.}\\
&\text{c}')\quad Z_3\ W\text{-asph.}
\end{aligned}
\]

donnent lieu aux implications

\[
a'\mathbin{\&}b'\Longrightarrow c',\qquad b'\mathbin{\&}c'\Longrightarrow a',
\]

et que si $W$ est stable par facteurs directs, alors

\[
a'\mathbin{\&}c'\Longrightarrow b'.
\]

C'est à peu de choses près le cor. de la page 76, complété à la page 77 (corollaire), à cela près qu'on y suppose les flèches des imm. ouvertes, et le carré cartésien \deleted{et \unsure{cocartésien}}\inserted{(\unsure{donc|et} aussi cocartésien)} dans $\mathrm{Cat}$, \uline{i.e.} $Z_3=Z_1\cup Z_2$, $Z_0\mathrel{\unsure{=|>}}Z_1\cap Z_2$. Mais
% archive sheet number: 317
